\chapter{多项式与交换域。}

交换域的理论——以及与它紧密相联系的多项式理论——直接来源于直到 19 世纪中叶一直构成经典代数之主要对象的东西：代数方程的求解，以及与之等价的几何作图问题。

一旦着手求解次数 \textgreater{} 1 的代数方程，人们就会遇到一类全新的计算困难：仅凭从给定量出发的“有理”计算，是不可能完成对一个未知数的确定的。这一困难必定在非常早的阶段就被察觉到了；而巴比伦人知道把二次方程和双二次方程的求解化归为一种新的代数运算——开平方——则必须算作他们最重要的贡献之一（流传到我们手中的文本里大量以这种方式求解的数值方程可以证明这一点（{[}232{]}，第~183-193~页））。就形式计算而言，古代在代数方程求解问题上从未超过这一点；古典时代的希腊人确实仅限于用几何的语言重新找到巴比伦人的公式，而以代数形式使用这些公式，在海伦（公元 100 年）和丢番图之前尚无明证。

希腊人是在一个完全不同的方向上取得决定性进展的。关于巴比伦人如何设想和计算非平方整数的平方根，我们所知甚少；\(^{1}\) 在流传到我们手中的涉及这一问题的罕见文本里，他们似乎满足于相当粗糙的近似方法（{[}232{]}，第~33-38~页）。毕达哥拉斯学派曾经严格地确定了可公度量的概念，并赋予它一种近乎宗教的性质，故而不可能持这种观点；而把他们最终引向证明这个数是无理数的，可能正是把 \(\sqrt{2}\) 表为有理数的反复尝试的失败。\(^{2}\)

我们将在别处（参见第~148~页）叙述这一标志着数学史重大转折的发现如何深刻地反作用于希腊人的“数”的概念，并如何引导他们创立一种纯几何特征的代数，以便在其中找到一种表示不可公度之比（或者也许是为这种比给出一种“存在性”证明）的办法——他们拒绝把这些比当作数。在大多数情形，他们把一个代数问题化归为两条适当选取的辅助平面曲线的相交，或者化归为这类相交的若干次相继确定。一些迟出的、权威性甚小的传统把这类作图的一种最初的分类——它注定有长久而辉煌的历程——归于柏拉图：看来是出于比数学更偏哲学的理由，他把所谓“直尺和圆规”的作图，也就是说除直线和圆之外不用任何辅助曲线的那些作图，单独放到一边。\(^{3}\) 无论如何，欧几里得在其《原本》{[}107{]} 中仅限于专门处理可用这种方式求解的问题（然而并未以一个特别的名称来刻画它们）；这一情况无疑大大有助于把此后几个世纪数学家的注意力固定在这些问题上。但我们今天已经知道，\(^{4}\) 可以用“直尺和圆规”求解的代数方程属于一种非常特殊的类型；特别地，三次不可约方程（在有理数域上）不能以这种方式求解，而希腊人很早便遇到了一些一直享有盛名的三次问题，例如倍立方（方程 \(x^{3}=2\) 的解）和三等分角；化圆为方则相反使他们面对一个超越的问题。我们看到，为了解决这些问题，他们引入了许多代数的（圆锥曲线、狄奥克莱斯的蔓叶线、尼科墨德斯的蚌线）或超越的（狄诺斯特拉图的割圆曲线、阿基米德螺线）曲线；这不可能不把他们引向对这些不同曲线的自主研究，从而为解析几何、代数几何和无穷小演算未来的发展准备了道路。但这些方法对代数方程的求解没有带来任何进展，\(^{5}\) 而古代唯一对这一问题作出了像样贡献、并对中世纪和文艺复兴时期的代数学家施加了持久影响的著作，就是欧几里得《原本》{[}107{]} 的第 X 卷；在这一卷中（有些史家想把其中的主要结果归于泰阿泰德），他考

察由若干根式的组合所得的表达式，例如 \(\sqrt{\sqrt{a}} \pm \sqrt{b}\)（a、b 为有理数），给出这些表达式为无理数的条件，把它们分成许多类别（他证明这些类别互不相同），并研究这些各种各样的无理数之间的代数关系，例如我们今天会写成

\[
\sqrt {\sqrt {p} + \sqrt {q}} = \sqrt {1 / 2 (\sqrt {p} + \sqrt {p - q})} + \sqrt {1 / 2 (\sqrt {p} - \sqrt {p - q})};
\]

这一切都是用《原本》惯用的几何语言表述的，这使得阐述格外沉闷而笨拙。

在古典希腊数学衰落之后，有关代数方程的各种概念发生了变化。毫无疑问，在整个古典时期，希腊人掌握着平方根的不定近似方法，可惜我们对之了解得很差。\(^{6}\) 随着印度人，然后是中世纪的阿拉伯人及其西方效仿者，各次开方变成了一种倾向于被视为与代数的有理运算同等基本的运算，并且像后者一样，用越来越便于计算的记号来表示。\(^{7}\) 在整个中世纪中日臻完善的二次方程理论（根的个数、负根、不可能的情形、二重根）以及双二次方程理论，给出了方程“根式解”公式的范本，几个世纪里代数学家们都试图从中仿造出解更高次方程、首先是三次方程的类似公式。13 世纪把阿拉伯科学介绍到西方的主要人物比萨的莱昂纳多认识到，无论如何，欧几里得在其第 X 卷中分类的那些无理数不能为此目的所用（在这样一门包含如此多不可能性的理论中，这又是一个不可能性的证明），而且我们看到他已经在尝试用立方根做类似的计算，得到了诸如

\[
\sqrt [ 3 ]{16} + \sqrt [ 3 ]{54} = \sqrt [ 3 ]{250},
\]

这样的关系式，它们与欧几里得的平方根公式类似（此外，更早的例子在阿拉伯人那里已经可以找到）。但还要再经过三个世纪徒劳的努力，费罗才在 16 世纪初最终对方程 \(x^{3} + ax = b\) 得到了求解公式 \(^{8}\)

\[
x = \sqrt [ 3 ]{\frac {b}{2} + \sqrt {\left(\frac {b}{2}\right) ^ {2} + \left(\frac {a}{3}\right) ^ {3}}} + \sqrt [ 3 ]{\frac {b}{2} - \sqrt {\left(\frac {b}{2}\right) ^ {2} + \left(\frac {a}{3}\right) ^ {3}}}.\tag{5.1}
\]

我们在这里无法描述这一轰动性发现之富有色彩的一面——它一方面在塔尔塔利亚与另一方面在卡尔丹及其学派之间激起的争吵——也无法描述作为主角的学者们往往颇为动人的形象。但我们必须指出，随后在卡尔丹及其学生手中，方程理论取得了决定性的进步。卡尔丹对使用负数的反感比他大多数同时代人要少，他由此观察到三次方程可以有三个根，而双二次方程有四个根（{[}50{]}，第 IV 卷，第~259~页），并且他指出，\(x^3 + bx = ax^2 + c\)（他知道如何消去这个方程中的 \(x^2\) 项）的三个根之和总等于 \(a\)（出处同上）。无疑是在这一关系以及对其普遍性质的直觉的引导下，他萌生了根的重数这一概念的最初想法；尤其是，他壮起胆子（并非没有一番言辞上的铺垫）对含负数平方根的表达式进行形式计算。他很可能是被如下事实引到这一步的：当 \(\left(\frac{b}{2}\right)^{2}+\left(\frac{a}{3}\right)^{3}<0\) 时（这是一种被称为“不可约”的情形，卡尔丹已经认识到其中存在三个实根），这样的表达式在公式（1）的使用中会自然出现；无论如何，这一点在他的学生 R. 邦贝利那里清楚地显现出来，邦贝利在其《代数学》（{[}28 a{]}，第~293~页）中证明了关系式

\[
\sqrt [ 3 ]{2 + \sqrt {-121}} = 2 + \sqrt {-1}
\]

并且注意以与现代的阐述十分接近的形式明确给出复数的计算规则。\(^{9}\) 最后，早在 1545 年，卡尔丹的另一个学生 L. 费拉里就借助一个三次的辅助方程成功地解出了普遍的四次方程。\(^{10}\)

在如此迅速的进展之后，直到 18 世纪中叶的随后一个时期，几乎只是发展意大利学派引入的那些新观念。得益于韦达为代数记号带来的这一本质进步，他得以用一般的方式表达代数方程的系数与根之间的关系，至少当诸根皆为正数时是如此 \(^{11}\)（{[}319{]}，第~158~页）。更大胆的 A. 吉拉尔 {[}129{]} 则毫不犹豫地断言（当然没有证明）：只要把“不可能的根”各按其重数计入，n 次方程就恰好有 n 个根，而且这些根满足韦达给出的那些关系；他还第一次得到了根的幂和直到指数 4 的表达式。

但 17 世纪的精神转向了其他方向，代数只是作为间接的后果，才从解析几何和无穷小演算的新发现中得到些许好处。于是，与笛卡儿求代数曲线切线的方法（参见第~178~页）相联系，出现了代数方程之根的重数判别法，它由他的学生许德（{[}85 b{]}，第 I 卷，第~433~页及第~507-509~页）陈述。代数函数与超越函数之间的区分，无疑也须追溯到笛卡儿的影响，它平行于他在其《几何学》中就“几何”曲线与“机械”曲线所作的区分（参见第~176~页和第~193~页）。无论如何，这一区分在 J. 格雷戈里那里是完全清楚的，他在 1667 年甚至试图证明圆扇形的面积不可能是由弦和半径表出的代数函数（{[}136 a{]}；参见{[}155{]}）。“超越的”一词出自莱布尼茨，他在其整个生涯中从未间断对这类分类问题的兴趣，并在大约 1682 年发现了格雷戈里所寻求结果的一个简单证明，即证明了 sin x 不可能是 x 的代数函数（{[}198 a{]}，第 V 卷，第~97-98~页）。\(^{12}\) 他与他的朋友奇恩豪斯，是他那个时代仅有的几位仍对代数方程“根式解”问题感兴趣的数学家。在其早期阶段，我们看到他研究三次方程的“不可约情形”，并确信（事实上并无充分的证明）在这种情形下不可能从求解公式中除去虚量（{[}198 d{]}，第~547-564~页）。大约在同一时期，他也未能成功地攻克五次方程的根式解问题；后来，当奇恩豪斯声称借助形如 \(y = P(x)\) 的变换（其中 P 是一个适当选取的四次多项式）使方程中除两端两项外的所有项都消失、从而解决了问题时，莱布尼茨立即看出，确定 \(P(x)\) 的系数的那些方程的次数 \textgreater{} 5，并断定这个方法注定要失败（{[}198 d{]}，第~402-403~页）。

看来，正是新分析学的需要使人们对代数的兴趣逐渐恢复起来。由莱布尼茨和约翰·伯努利完成的有理函数的积分法，以及与之紧密相联系的虚对数问题，成为深化虚数运算的契机，也使多项式分解为一次因子的这一问题（“代数基本定理”）得到回答。\(^{13}\) 从 18 世纪初起，二项方程 \(x^{n}-1=0\) 的求解就被科茨和棣莫弗化归为把圆分成 n 等份；因此，要得到其根的“根式”表达式，只需会做 n 为奇素数的情形，而棣莫弗注意到，代换 \(y = x + \frac{1}{x}\) 那时便把问题化为解一个 \((n - 1)/2\) 次方程的“根式解”。至于“基本定理”，在一般解法的“根式解”屡次受挫（包括欧拉的若干次尝试（{[}108 a{]}，（1），第 VI 卷，第~1-19~页和第~170-196~页））之后，人们开始寻求先验的、不使用显式求解公式的证明。不必深入所提诸方法的细节（它们最终引出拉格朗日和高斯的证明；参见第~91~页和第~161~页），这里应当指出 18 世纪中叶人们看待这一问题所取的观点：人们假定（除了对一般性的一种模糊感觉——它无疑与 A. 吉拉尔那里一样，来自系数与根之间存在关系这一事实——之外，没有任何别的理由）n 次方程总有 n 个“理想”根，可以对它们像对数那样进行计算，却不知道它们是否为（实的或复的）数；而必须证明的则是（必要时可使用关于理想根的计算）：这些根中至少有一个是通常的复数。\(^{14}\) 以这种有缺陷的形式，人们认识到这里已经有了“形式添加”这一一般观念的第一个胚胎，它不顾高斯的反对（{[}124 a{]}，第 III 卷，第~1~页），将成为现代交换域理论的基础。

随着拉格朗日（{[}191{]}，第 III 卷，第~205-421~页）和范德蒙德 {[}315{]} 的那些基本论文，1770 年开启了代数方程理论历史上一个新的决定性时期。继直到那时一直无可争议地占统治地位的、多少有些碰运气的求解公式尝试的经验主义之后，随之而来的是对所述问题以及可望解决它们的方法的系统分析，这一分析在六十年后将引出伽罗瓦的最终结果。拉格朗日和范德蒙德都从根式的多个取值在次数 \(\leq 4\) 的方程求解公式中引入的多义性出发；这一事实曾引起欧拉的注意（{[}108 a{]}，（1），第 VI 卷，第~1-19~页），他特别指出了在费罗的公式中应当如何把其中出现的各根式的取值适当配对，以便得到 3 个根而不是 9 个。\(^{15}\) 拉格朗日注意到，费罗公式中的每个立方根都可以写成 \(\frac{1}{3}(x_1 + \omega x_2 + \omega^2 x_3)\) 的形式，其中 \(\omega\) 是 1 的一个立方根，\(x_1, x_2, x_3\) 是所论方程的三个根，按某种次序取出；他作出了基本的观察：三个根的函数 \((x_1 + \omega x_2 + \omega^2 x_3)\) 在三个根的所有置换下只能取两个不同的值，这就先验地解释了求解这个方程的那些方法何以成功。对四次方程求解方法的类似分析把他引到四个根的函数 \(x_{1}x_{2} + x_{3}x_{4}\)，它在根的所有置换下只取三个不同的值，因而是以所给方程的系数的有理函数为系数的一个三次方程的根；\(^{16}\) 拉格朗日说，这些事实构成了“真正的原则，以及所谓三次和四次方程求解的形而上学 \(^{17}\)”（{[}191{]}，第 III 卷，第~357~页）。凭借这些例子，他提议一般地对 n 次方程研究有理函数 V 在 n 个根被任意置换时所能取的值的个数 \(\nu\)；\(^{18}\) 实际上，他由此开创了（用的是仍紧密依附于方程论的术语）群的理论和域的理论，并且已经用与今天所用的相同的原则得到了它们的几个基本结果。例如，他借助今天用来证明有限群的子群的阶整除群的阶的那一论证，证明了数 \(\nu\) 是 n! 的一个因子。更为引人注目的是这样一个定理：若 \(V_{1}\) 和 \(V_{2}\) 是根的两个有理函数，且 \(V_{1}\) 与 \(V_{2}\) 在同样的置换下不变，则其中每一个都是另一个与方程的系数的有理函数（这是伽罗瓦把伽罗瓦扩张的子扩张刻画为其伽罗瓦群的不变量域的定理的一个特殊情形）：他说，“这个问题在我看来是方程论中最重要的问题之一，我们将给出的它的一般解，将是使代数的这一部分焕然一新的手段”（{[}191{]}，第 III 卷，第~374~页）。

所有这些研究在拉格朗日的心目中自然只是准备步骤，为的是分析通过逐次化归为次数较低的方程来求解代数方程的可能方法，而他证明了这样的方法联系于构造在根的置换下取值少于 n 个的根的有理函数。无疑是在他关于三次方程的结果的引导下，他一般地引入了“拉格朗日预解式”\(y_{k} = \sum_{h=1}^{n} \omega_{k}^{h} x_{h}\)，其中 \(\omega_{k}\) 是一个 n 次单位根（\(1 \leq k \leq n\)），他清楚地表明这 n 个数的知识如何导出诸根 \(x_{k}\) 的知识，并一般地寻求 \(y_{k}\) 所满足的方程的次数；例如他证明，若 n 为素数，则诸 \(y_{k}\) 是一个 n-1 次方程的根，其系数是某个 \((n-2)!\) 次方程的一个根的有理函数，而该方程的系数又可由所给方程的系数有理地表出。他总结道：“如果我没有弄错的话，这些就是方程求解的真正原则，以及达到这一目标的最正确的分析；正如可以看到的，一切都可以归结为一种组合计算，借助它，人们可以先验地得到所应期待的结果”（{[}191{]}，第 III 卷，第~403~页）。

至于范德蒙德的论文，它独立于拉格朗日的论文，却在许多点上与后者相遇，特别是在寻求在根的置换下取尽可能少的不同值的根的有理函数这一观念方面，\(^{19}\) 以及他为此目的同样引入的“拉格朗日预解式”的研究方面。他的著作远不具有拉格朗日著作的那种清晰性和普遍性；然而在一点上它明显走得更远，即把同样的观念应用于 n 为奇素数的分圆方程 \(x^{n}-1=0\)。拉格朗日满足于提醒读者，这个方程可化归为一个 \(m=(n-1)/2\) 次的、以有理数为系数的方程，而当 \(n\geq11\) 时并不去求解它；范德蒙德则断言，由于 \(x^{n}-1=0\) 的各个根之间的关系，这个方程的拉格朗日预解式的 m 次幂是有理数；但他仅限于在 n=11 的情形核验这一断言的真实性，而未以一般的方式证明它。

只是到了 30 年后，范德蒙德所预告的结果才由 C. F. 高斯完全证明。\(^{20}\) 他关于方程 \(x^n - 1 = 0\)（\(n\) 为奇素数）的决定性结果，属于他那令人难忘的算术研究的总纲领（{[}124 a{]}，第 I 卷，第~413~页以下），并以一种非常特别的方式显示出他驾驭我们今天所谓的循环群理论的娴熟技艺。在证明了多项式 \(\Phi_n(x) = (x^n - 1)/(x - 1)\) 对奇素数 \(n\) 不可约 \(^{21}\) 之后，他想到把它的 \(n - 1\) 个根写成 \(\zeta^{g^k} = \zeta_k (0 \leq k \leq n - 2)\) 的形式，其中 \(g\) 是同余式 \(z^{n-1} \equiv 1 \pmod{n}\) 的一个原根）（用现代的语言说，这相当于挑明 \(\Gamma\) 即方程 \(\Phi_n(x) = 0\) 的群是循环群这一事实）。对每一个因子 \(e\)（即 \(n-1\) 的因子），他对应地取 \(f = (n-1)/e\) 个“周期”\(\eta_\nu = \xi_\nu + \xi_{\nu+f} + \xi_{\nu+2f} + \cdots + \xi_{\nu+(e-1)f} (1 \leq \nu \leq f)\)，并实质上证明了：诸 \(\eta_{\nu}\) 的以有理数为系数的线性组合构成一个域，它由 \(f\) 个周期 \(\eta_{\nu}\) 中的任何一个生成，并且它在有理数域上的次数为 \(f\)（这个域自然对应于 \(\Gamma\) 的 \(e\) 阶子群）。我们在这里无法深入他的分析的细节以及由此得出的重要算术推论；只需注意到，它特别使他得到了那个关于可以用“直尺和圆规”作出边数等于形如 \(2^{2^k} + 1\) 的素数的多边形的著名定理。至于方程 \(\Phi_n(x) = 0\) 的根式解，它容易由把周期理论应用于一个拉格朗日预解式的 \(f\) 次幂（预解式为 \(\sum_{\nu=0}^{f-1} \omega^\nu \eta_\nu\)，其中 \(\omega^f = 1\)）而得出。

他的同胞鲁菲尼的研究 {[}265{]} 与《算术研究》同时代，而且直接承自拉格朗日；他在拉格朗日停下的地方重新拾起这一问题，提出的目标是证明“一般”五次方程 \(^{24}\) 的“根式解”之不可能性。鲁菲尼的证明冗长而晦涩，尽管几经修改仍不完全；但它已经非常接近阿贝尔后来所得到的（原则上正确的）证明。\(^{25}\) 它的主要意义尤其在于引入了置换的计算以及群论的最初概念，鲁菲尼发展这些内容是为了证明：不存在方程 5 个根的一个函数，当诸根被任意置换时它取的值多于 2 个而少于 5 个。

我们已经说过（见第~51~页），置换群理论的这一最初草稿是如何在几年之后由柯西发展和系统化的。但是，如果说就置换而言、拉格朗日诸观念的发展所必需的那些概念就这样被逐渐澄清了，那么还必须以同样清晰的方式陈述域论的最初原理。这正是鲁菲尼所缺少的东西，而这也正是阿贝尔和伽罗瓦在代数方程求解问题的最后阶段将要做到的。

在他短促的一生中，阿贝尔从未间断过对这一问题的关切。当他还几乎是个孩子的时候，他曾相信自己已经得到了一般五次方程的根式求解公式。后来认识到了自己的错误，他便不达证明出这样的公式并不存在的目的誓不罢休（{[}1{]}，第 I 卷，第~66~页）。但他并未止步于此。当他的竞争者雅可比作为分析学家发展椭圆函数理论时，支配着阿贝尔这方面工作的却是代数的观点，其中心是椭圆函数的分（割）方程理论（{[}1{]}，第 I 卷，第~265、377~页及各处）。他由此用仿照高斯解分圆方程的方法（{[}1{]}，第 I 卷，第~310~页和第~358~页）得到了一些可用根式求解的新类型的方程；\(^{26}\) 由这一结果他进而上升到“阿贝尔”方程的概念，并在一篇著名论文中证明了这类方程可用根式求解（{[}1{]}，第 I 卷，第~478~页）；也正是在这个场合，他精确地定义了在一个给定域（由他所研究的方程的系数生成的域）上不可约的多项式的概念。\(^{27}\) 最后，死亡于 1829 年把他击倒了，那时他正好着手刻划所有可用根式求解的方程的一般问题，并且刚刚向克雷勒和勒让德通报了已经非常接近伽罗瓦的结果（{[}1{]}，第 II 卷，第~219-243、269-270~页和第~279~页）。

三年之后，为这座大厦加冕的任务留给了他 {[}123{]}。像阿贝尔一样，但以一种更为清晰的方式，他一开始就（除术语差异外）定义了一个数属于由给定的诸数所生成的域、添加的概念，以及给定域上的不可约多项式。给定一个方程 \(F(x) = 0\)，它没有重根、系数属于一个给定的域 \(K\)，他相继证明：“总可以构造诸根的一个函数 \(V\)，使得以一切可能的方式在这个函数中置换诸根时所得到的诸值互不相等”，这个函数“具有这样的性质：所提方程的所有根都可以表示为 \(V\) 的有理函数”，以及，设 \(V, V', V'', \ldots\) 是 \(V\) 所满足的不可约方程的诸根，“若 \(a = f(V)\) 是所提方程的一个根，则 \(f(V')\) 同样将是所提方程的一个根”（{[}123{]}，第~47-51~页）；用现代的语言说，他由此证明了 \(V\) 以及它在 \(K\) 上的任一共轭元都生成域 \(N\)，即 \(F\) 的诸根所构成的域。然后他把 \(\Gamma\) 定义为 \(F\) 的群，即诸根 \(x_i\) 的如下置换的集合：把 \(V\)——在各个 \(x_i\) 作为 \(V\) 的函数的有理表达式中——代之以 \(V\) 的任一共轭元，便得到这些置换；他并由此立即得到 \(K\) 中元素凭其在 \(\Gamma\) 的每个置换下不变这一性质的基本刻划（{[}123{]}，第~51~页）。接着他证明，若 \(N\) 包含另一个多项式的根域 \(L\)，则 \(N\) 在 \(L\) 上的群是 \(\Gamma\) 的一个正规子群（这个概念正是他在此时引入的）（{[}123{]}，第~175~页）。最后他由此导出方程可用根式求解的判据，其论证的本质部分如下：设基域 \(K\) 包含所有的单位根，则依假设必存在一个递增序列 \((K_i)_{0 \leq i \leq m}\)，它由介于 \(K\) 与 \(N\) 之间的中间域组成，满足 \(K_0 = K, K_m = N\)，域 \(K_{i+1}\) 是通过向 \(K_i\) 添加一个二项方程 \(x^{n_i} - a_i = 0\)（其中 \(a_i \in K_i\)）的所有根而得到的。因此在 \(\Gamma\) 中应存在一个子群的递减序列 \((\Gamma_i)\)，使得 \(\Gamma_0 = \Gamma, \Gamma_m = \varepsilon\)（中性元），\(\Gamma_{i+1}\) 在 \(\Gamma_i\) 中正规，且商群 \(\Gamma_i / \Gamma_{i+1}\) 是循环群（这时就说群 \(\Gamma\) 是可解的）。反之，若情况如此，拉格朗日预解式的使用表明 \(K_{i+1}\) 是由向 \(K_i\) 添加某个二项方程的所有根而得到的，从而方程 \(F(x) = 0\) 可用根式求解。\(^{28}\) 于是，“一般”方程——其次数 \(n > 4\)——根式解之不可能，便是如下事实的推论：这个方程的群 \(\Gamma\) 同构于对称群 \(S_n\)，而后者不是可解的。

从 19 世纪中叶起，如我们已经指出的那样（参见第~53~页），代数学家们大大扩展了他们的研究范围，而在此之前，这一范围几乎完全局限于方程的研究。在伽罗瓦发现的启示下，人们看出“根式解”问题只不过是无理数分类这一一般问题的一个相当人为的特殊情形。在整个 19 世纪后半叶，人们将从多个方面进攻的正是后一个问题，许多零散的结果将逐渐积累起来，为施泰尼茨的综合准备了道路。

首先就有代数无理数而言，分类的一条基本原理由伽罗瓦理论提供，它把代数方程的研究化归为其群的研究。事实上，在纯粹代数中，此时研究的主要对象是置换群的理论，对此我们在这里无须谈论（参见第~53~页）。代数数域理论的其他进展来自同一时期数论和代数几何的发展。这些进展主要涉及阐述理论的方式，而且其中大部分应归功于戴德金 {[}79{]}，他引入了域和环的概念，\(^{29}\) 并（联系到他对超复系统的研究）系统地发展了扩张理论的线性方面（{[}79{]}，第 III 卷；第~33~页以下）。也是他把伽罗瓦群视为由所论扩张的自同构所组成，而不再仅仅视为方程诸根的一个置换群；他（就数域）证明了关于自同构线性无关的基本定理（{[}79{]}，第 III 卷，第~29~页）以及伽罗瓦扩张的正规基的存在性（{[}79{]}，第 II 卷，第~433~页）。最后，他着手处理无限次代数扩张的问题，并注意到伽罗瓦理论不能原封不动地应用于它（伽罗瓦群的任意子群并不总是与一个扩张关于某个子扩张的群相同）；而且凭着一种大胆的直觉，他已经开始把伽罗瓦群设想为一个拓扑群 \(^{30}\)——这一思想直到 1928 年才由克鲁尔发展的无限次伽罗瓦扩张理论开花结果 {[}187 d{]}。

与这一演变并行，域上的超越元素的概念变得精确了。超越数的存在性第一次由刘维尔于 1844 年借助一个基于丢番图逼近理论的显式构造得到证明（{[}204 c{]}）；康托尔于 1874 年用关于集合的势的简单考虑给出了另一个“非构造性”证明（{[}47{]}）；最后，埃尔米特于 1873 年证明了 e 的超越性，林德曼于 1882 年用与埃尔米特类似的方法证明了 \(\pi\) 的超越性，从而给化圆为方这一古老问题盖上了最终的封印。\(^{31}\)

至于超越数在代数计算中所起的作用，克罗内克于 1882 年注意到，若 x 在域 K 上是超越的，则域 \(K(x)\) 同构于有理分式域 \(K(X)\)（{[}186 a{]}，第 II 卷，第~253~页）。他还把未定元向域的添加作为其代数数理论阐述的基石（{[}186 a{]}，第 II 卷，第~245-387~页）。另一方面，戴德金和韦伯于同年 {[}80{]} 表明，算术的方法如何能帮助建立代数曲线的理论。由此可以看到算术与代数几何之间的类似在几个方向上显现出来，而这些类似对双方都将显示出极大的丰饶。

在所有这些研究中，出现的域都是由经典数学意义下的“具体”元素——（复）数或复变量的函数——组成的。\(^{32}\) 但早在 1882 年，克罗内克就充分认识到（高斯和伽罗瓦都曾隐约感觉到）这样一个事实：“未定元”在他的理论中所起的作用只是某个代数的基元素的作用，而不是分析意义下的变量的作用（{[}186 a{]}，第 II 卷，第~339-340~页）；1887 年，他发展了这一思想，并将其同一个庞大的纲领联系在一起，该纲领的目标不外乎重新铸造全部数学，摒弃一切不能化归为整数上的代数运算的东西。正是在这个场合，克罗内克接过柯西（{[}56 a{]}，（1），第 X 卷，第~312~页和第~351~页）的一个想法——柯西曾把复数域 C 定义为剩余域 \(\mathbf{R}[X]/(X^{2}+1)\)——并表明代数数理论完全独立于“代数基本定理”，甚至独立于实数理论：所有（有限次的）代数数域都同构于一个剩余域 \(\mathbf{Q}[X]/(f)\)（f 是 \(\mathbf{Q}\) 上的不可约多项式）（{[}186 a{]}，第 III\(_{1}\) 卷，第~211-240~页）。几年之后，H. 韦伯 {[}327 a{]} 在发展域的公理化理论的一份最初草稿时指出，克罗内克的这一方法实际上适用于一切基域 K。韦伯特别指出，可以取 \(\mathbf{Z}/(p)\)（p 为素数）作为 K，从而使“模 p”同余的演算回归到域论；后者诞生于 18 世纪下半叶，与欧拉、拉格朗日、勒让德和高斯的名字相联，而它所显示出的与代数方程理论的类似并未逃过观察者的眼睛；发展这一类似，伽罗瓦（心中惦记着他关于群论的研究）毫不犹豫地引入了模 p 不可约同余的“理想根”，\(^{33}\) 并指出了其主要性质（{[}123{]}，第~113-127~页）。\(^{34}\) 当把克罗内克的方法应用于 \(\mathbf{Z}/(p)\) 时，重新得到的（除术语外）乃是塞尔和戴德金（{[}79{]}，第 I 卷，第~40~页）已经给出的关于这些“伽罗瓦虚数”的理论的表述。

在世纪之交，所有这些“抽象域”的例子之外又添上了若干类型截然不同的新域，即韦罗内塞 {[}318{]} 引入的形式幂级数域，尤其是亨泽尔 {[}157 f{]} 的 p 进域。正是后者的发现促使施泰尼茨（正如他自己明确说明的那样）提炼出所有这些理论共有的抽象观念；他在一部奠基性的著作 {[}294 a{]} 中完成了这项工作，这部著作可以说孕育了代数的现今概念。他系统地推演交换域公理的各种推论，从而引入了素域、可分（代数）元、完备域等概念，定义了扩张的超越次数，并最终证明了任意域的代数闭扩张的存在。

近来，施泰尼茨的理论在若干重要方面得到了补充。一方面，阿廷的工作使伽罗瓦理论的线性特征凸显出来 {[}7 a{]}。另一方面，导子这一普遍概念（仿照经典微分学的形式性质而得）——戴德金（{[}79{]}，第 II 卷，第~412~页）已隐约感觉到它，施泰尼茨在有理分式域的特殊情形下引入了它（{[}294 a{]}，第~209-212~页）——在超越扩张的研究（对现代代数几何至关重要）中得到了成功的运用，尤其是运用于把可分性概念推广到这些扩张 {[}330 d{]}。

